\documentclass[../DoAn.tex]{subfiles}
\begin{document}

Chương trước đã trình bày ngữ cảnh bài toán, các nghiên cứu liên quan và các cơ sở lý thuyết phục vụ cho bài toán phát hiện và hiệu chỉnh lỗi ước lượng tư thế người 3D. Trên cơ sở đó, chương này trình bày phương pháp đề xuất theo hướng hậu xử lý kết quả từ mô hình monocular pose estimation, một hướng phù hợp với các nghiên cứu sử dụng mô hình có sẵn như WHAM, SMPL, SMPLify và Learnable SMPLify để tạo hoặc tinh chỉnh chuyển động người 3D \cite{shin2024wham,smpl2015,smplify2016,yang2025learnable}. Nội dung chương gồm tổng quan pipeline, mô hình hóa bài toán, định dạng dữ liệu qua từng mô-đun, các bước phát hiện lỗi, phân xử mâu thuẫn, hiệu chỉnh pose và thiết kế đánh giá thực nghiệm.

\section{Tổng quan giải pháp}
\label{sec:method_overview}

Phương pháp đề xuất được xây dựng cho bài toán hậu xử lý hai chuỗi pose 3D được ước lượng độc lập từ hai video monocular. Đầu vào ban đầu của hệ thống gồm hai video RGB ghi cùng một chuyển động từ hai góc nhìn khác nhau. Mỗi video được xử lý bởi một mô hình monocular, chẳng hạn WHAM, để sinh ra chuỗi pose 3D, tham số cơ thể SMPL, vector tịnh tiến và các biểu diễn liên quan \cite{shin2024wham,smpl2015}. Đồ án không thay thế mô hình WHAM, mà sử dụng đầu ra của mô hình này làm dữ liệu đầu vào cho pipeline phát hiện và hiệu chỉnh lỗi \cite{vu2026multiperspective}.

Mục tiêu của phương pháp là tạo ra một chuỗi pose 3D sau hiệu chỉnh có tính nhất quán cao hơn giữa hai góc nhìn, ổn định hơn theo thời gian và phù hợp hơn với ràng buộc động học cơ thể. Cách tiếp cận này dựa trên nhận định từ các nghiên cứu multi-view rằng thông tin từ nhiều góc nhìn có thể hỗ trợ giảm nhập nhằng chiều sâu trong 3D human pose estimation \cite{canonpose2021,mtftransformer2023,gu2024multiview}. Tuy nhiên, khác với các hệ thống multi-view trực tiếp ước lượng pose từ ảnh đa camera, phương pháp đề xuất chỉ xử lý hai chuỗi pose đã được sinh ra độc lập từ hai mô hình monocular \cite{vu2026multiperspective}.

Luồng xử lý tổng thể gồm sáu mô-đun chính. Mô-đun thứ nhất nạp dữ liệu từ hai video và các tệp output của mô hình monocular. Mô-đun thứ hai đồng bộ hai chuỗi pose theo thời gian, vì so sánh liên góc nhìn chỉ có ý nghĩa khi hai frame biểu diễn cùng một thời điểm chuyển động. Mô-đun thứ ba căn chỉnh hai chuỗi pose về cùng một không gian so sánh, tương tự vai trò của các kỹ thuật căn chỉnh hình học như Procrustes alignment trong các bài toán đa góc nhìn \cite{gu2024multiview}. Mô-đun thứ tư kiểm tra nhất quán đa góc nhìn để phát hiện khớp hoặc frame nghi ngờ. Mô-đun thứ năm phân xử và hợp nhất thông tin giữa hai góc nhìn. Mô-đun cuối cùng hiệu chỉnh pose bằng ràng buộc thời gian, ràng buộc động học và biểu diễn SMPL/SMPLify hoặc Learnable SMPLify \cite{smpl2015,smplify2016,yang2025learnable}.

Bảng~\ref{tab:module_io_overview} tóm tắt luồng xử lý theo từng giai đoạn để làm rõ pipeline bắt đầu từ video RGB và kết thúc ở chuỗi pose đã hiệu chỉnh.

\begin{table}[H]
\centering
\caption{Đầu vào và đầu ra của các mô-đun trong pipeline đề xuất}
\label{tab:module_io_overview}
\begin{tabular}{|p{3.1cm}|p{3.7cm}|p{4.2cm}|p{4.0cm}|}
\hline
\textbf{Giai đoạn} & \textbf{Đầu vào} & \textbf{Xử lý chính} & \textbf{Đầu ra} \\
\hline
Sinh pose ban đầu & Hai video RGB cùng ghi một chuyển động & Chạy WHAM hoặc mô hình monocular tương đương \cite{shin2024wham} & Hai chuỗi pose thô từ hai góc nhìn \\
\hline
Chuẩn hóa dữ liệu & Output WHAM gồm joints, pose, shape, trans, metadata & Kiểm tra định dạng, số frame, đơn vị và tên trường & Hai pose sequence có cấu trúc thống nhất \\
\hline
Đồng bộ thời gian & Hai chuỗi pose có thể lệch frame & Ghép frame theo timestamp, chỉ số frame hoặc offset chuyển động & Hai chuỗi pose tương ứng theo thời gian \\
\hline
Căn chỉnh không gian & Hai chuỗi pose đã đồng bộ & Đưa hai chuỗi về cùng hệ tọa độ so sánh \cite{gu2024multiview} & Hai chuỗi pose đã căn chỉnh \\
\hline
Kiểm tra nhất quán & Hai chuỗi pose đã căn chỉnh & Tính sai khác khớp, xương, vận tốc, gia tốc \cite{manipose2024} & Ma trận sai khác, mask lỗi, điểm tin cậy \\
\hline
Phân xử và hợp nhất & Hai chuỗi pose, mask lỗi, điểm tin cậy & Chọn view đáng tin cậy hoặc hợp nhất có trọng số \cite{vu2026multiperspective} & Chuỗi pose hợp nhất \\
\hline
Hiệu chỉnh cuối & Pose hợp nhất và các ràng buộc thời gian/động học & Làm mượt, duy trì chiều dài xương, tinh chỉnh SMPL nếu dùng \cite{smplify2016,yang2025learnable,physcap2020} & Chuỗi pose 3D sau hiệu chỉnh \\
\hline
\end{tabular}
\end{table}

\begin{figure}[htbp]
    \centering
    \includegraphics[width=\textwidth]{figures/kien_truc_tong_quan.png}
    \caption{Kiến trúc tổng quan của hệ thống đề xuất}
    \label{fig:proposed_method_overview}
\end{figure}

Hình~\ref{fig:two_view_logic} mô tả chi tiết cách hệ thống tiếp nhận hai view đầu vào, đưa chúng qua bước judgment để chọn nhánh xử lý phù hợp, rồi chuyển sang solver cơ sinh học hoặc refinement SMPL khi cần. Hình này làm rõ rằng pipeline không chỉ có một luồng xử lý tuyến tính, mà còn có cơ chế phân xử cục bộ giữa hai đầu vào monocular trước khi hiệu chỉnh cuối cùng.

\begin{figure}[htbp]
    \centering
    \includegraphics[width=\textwidth]{figures/Logic.png}
    \caption{Logic xử lý hai view đầu vào trong hệ thống đề xuất}
    \label{fig:two_view_logic}
\end{figure}

Từ góc nhìn phương pháp, đồ án đặt ở vị trí hậu xử lý hai chuỗi pose đã có, chứ không huấn luyện thêm một mô hình multi-view mới. Vì vậy, phần tiếp theo mô tả chi tiết định dạng dữ liệu và vai trò của từng mô-đun trong pipeline để phục vụ cài đặt và đánh giá thực nghiệm.

\section{Mô hình hóa bài toán và biểu diễn dữ liệu}
\label{sec:problem_formulation}

Giả sử hệ thống có hai video monocular quan sát cùng một người thực hiện cùng một chuyển động. Hai video này được ký hiệu là $V^{(1)}$ và $V^{(2)}$, tương ứng với hai góc nhìn khác nhau. Sau khi xử lý bằng WHAM hoặc mô hình monocular tương đương, hệ thống thu được hai chuỗi pose 3D:
\[
X^{(1)} = \{X^{(1)}_t\}_{t=1}^{T_1}, \quad
X^{(2)} = \{X^{(2)}_t\}_{t=1}^{T_2}
\]
trong đó $X^{(v)}_t \in \mathbb{R}^{J \times 3}$ là tập tọa độ 3D của $J$ khớp tại frame $t$ từ góc nhìn $v$. Cách biểu diễn pose bằng tập keypoint 3D là biểu diễn phổ biến trong 3D human pose estimation và cũng là cơ sở cho các độ đo như MPJPE hoặc PA-MPJPE \cite{manipose2024}.

Ngoài tọa độ khớp 3D, mỗi frame có thể chứa thêm các biểu diễn SMPL gồm tham số pose $\theta_t$, tham số shape $\beta_t$, vector tịnh tiến $\tau_t$ và tập đỉnh mesh $M_t$. SMPL biểu diễn cơ thể người bằng lưới tham số, trong đó shape điều khiển hình dạng cơ thể và pose điều khiển biến dạng tư thế của mô hình \cite{smpl2015}. Việc giữ đồng thời joint-based representation và SMPL-based representation giúp pipeline vừa đánh giá được sai khác giữa các khớp, vừa có cơ sở áp dụng ràng buộc hình học và động học \cite{smpl2015,smplify2016}.

Bài toán của đồ án được mô hình hóa như sau. Với hai chuỗi pose monocular $X^{(1)}$ và $X^{(2)}$, cần tìm một chuỗi pose hiệu chỉnh $\hat{X}$ sao cho $\hat{X}$ thỏa mãn ba điều kiện. Thứ nhất, $\hat{X}$ nhất quán hơn với các vùng pose đáng tin cậy từ hai góc nhìn, dựa trên nguyên tắc multi-view consistency \cite{canonpose2021,mtftransformer2023,vu2026multiperspective}. Thứ hai, $\hat{X}$ ổn định theo thời gian và hạn chế các biến đổi đột ngột không hợp lý, tương ứng với vấn đề temporal stability trong monocular motion capture \cite{shin2024wham,physcap2020}. Thứ ba, $\hat{X}$ duy trì các ràng buộc động học cơ bản như chiều dài xương và quan hệ giữa các khớp, phù hợp với các nghiên cứu về mô hình cơ thể và cơ sinh học \cite{keller2023skin,delp2007opensim}.

Bài toán này có thể được mô tả dưới dạng tối ưu hậu xử lý:
\[
\hat{X} = \arg\min_X
\lambda_v L_{\text{view}}(X)
+ \lambda_t L_{\text{temp}}(X)
+ \lambda_b L_{\text{bone}}(X)
+ \lambda_s L_{\text{smpl}}(X)
\]
trong đó $L_{\text{view}}$ đo mức mâu thuẫn giữa hai góc nhìn, $L_{\text{temp}}$ đo độ không ổn định theo thời gian, $L_{\text{bone}}$ đo sai lệch ràng buộc skeleton, và $L_{\text{smpl}}$ đo mức phù hợp với biểu diễn SMPL. Các nhóm ràng buộc này kế thừa ý tưởng từ multi-view consistency, temporal refinement, skeleton constraint và SMPL-based refinement trong các nghiên cứu trước \cite{canonpose2021,physcap2020,smplify2016,yang2025learnable}. Các hệ số $\lambda_v$, $\lambda_t$, $\lambda_b$ và $\lambda_s$ điều khiển vai trò tương đối của từng nhóm ràng buộc. Công thức này đóng vai trò mô tả định hướng lý thuyết; lựa chọn cụ thể của các hàm loss sẽ được xác định trong phần cài đặt và thực nghiệm.

Bảng~\ref{tab:data_format} mô tả các trường dữ liệu chính được sử dụng trong pipeline. Các tên trường có thể thay đổi theo định dạng output thực tế của mô hình monocular, nhưng về mặt phương pháp cần bảo toàn các nhóm thông tin tương ứng.

\begin{table}[H]
\centering
\caption{Các trường dữ liệu trung gian trong pipeline}
\label{tab:data_format}
\small
\begin{tabular}{|p{3.0cm}|p{2.3cm}|p{2.4cm}|p{3.2cm}|p{4.2cm}|}
\hline
\textbf{Trường dữ liệu} & \textbf{Kích thước tham khảo} & \textbf{Bắt buộc/tùy chọn} & \textbf{Vai trò trong pipeline} & \textbf{Ghi chú} \\
\hline
\texttt{joints3d} & $T \times J \times 3$ & Bắt buộc & Tọa độ khớp 3D dùng cho căn chỉnh, so sánh và đánh giá & Cơ sở cho các bước alignment, consistency và refinement \\
\hline
\texttt{joints2d} & $T \times J \times 2$ & Tùy chọn & Dùng cho chiếu lại lên ảnh hoặc kiểm tra trực quan & Hữu ích khi cần đối chiếu với video gốc \\
\hline
\texttt{pose} & $T \times P$ & Tùy chọn nếu chỉ xử lý keypoint; cần nếu dùng SMPL & Tham số tư thế trong mô hình cơ thể & Cần đồng bộ với SMPL/learned refinement \\
\hline
\texttt{betas} / \texttt{shape} & $B$ hoặc $T \times B$ & Tùy chọn & Tham số hình dạng cơ thể trong SMPL & Có thể giữ cố định theo người hoặc theo frame \\
\hline
\texttt{trans} & $T \times 3$ & Tùy chọn nhưng nên có & Vector tịnh tiến toàn cục & Quan trọng cho hệ tọa độ thế giới và đánh giá chuyển động \\
\hline
\texttt{vertices} & $T \times V \times 3$ & Tùy chọn & Mesh 3D; chỉ cần nếu refinement trên SMPL/mesh & Tăng chi phí nhưng hỗ trợ kiểm tra hình học \\
\hline
\texttt{confidence} & $T \times J$ & Tùy chọn & Độ tin cậy từng khớp, dùng trong phân xử và fusion & Có thể đến từ mô hình monocular hoặc hậu xử lý \\
\hline
\texttt{aligned\_joints3d} & $T \times J \times 3$ & Trung gian & Pose sau khi đưa về cùng không gian so sánh & Dùng cho consistency check \\
\hline
\texttt{consistency\_score} & $T \times J$ & Trung gian & Điểm sai khác giữa hai view & Dùng để đánh dấu khớp/frame nghi ngờ \\
\hline
\texttt{error\_mask} & $T \times J$ & Trung gian & Đánh dấu khớp/frame nghi ngờ lỗi & Làm đầu vào cho judgment/fusion \\
\hline
\texttt{fused\_joints3d} & $T \times J \times 3$ & Trung gian & Pose sau phân xử và hợp nhất & Là đầu ra của bước fusion trước refinement \\
\hline
\texttt{refined\_joints3d} & $T \times J \times 3$ & Đầu ra & Pose cuối cùng sau hiệu chỉnh & Dùng cho đánh giá và trực quan hóa \\
\hline
\end{tabular}
\end{table}

\section{Module 1: Nạp dữ liệu và kiểm tra đầu vào}
\label{sec:data_loader}

Mô-đun nạp dữ liệu có nhiệm vụ đọc hai video RGB và hai tệp output từ mô hình monocular. Trong phạm vi đồ án, output từ WHAM được xem là nguồn dữ liệu chính vì WHAM cung cấp các biểu diễn phục vụ tái tạo chuyển động 3D trong hệ toàn cục \cite{shin2024wham}. Mô-đun này không thay đổi pose, mà kiểm tra dữ liệu đầu vào và chuyển đổi về một định dạng thống nhất cho các bước sau.

Đầu vào của mô-đun gồm đường dẫn video, đường dẫn tệp output và metadata như số frame, tốc độ khung hình, kích thước ảnh và thông số camera nếu có. Metadata video được dùng để kiểm tra tính đồng bộ và tính đầy đủ của dữ liệu. Đầu ra là hai đối tượng dữ liệu pose sequence, mỗi đối tượng chứa chuỗi khớp 3D, tham số SMPL, vector tịnh tiến và các thông tin phụ trợ \cite{shin2024wham,smpl2015}.

Quy trình xử lý của mô-đun này gồm các bước: đọc metadata của hai video, đọc output monocular tương ứng, kiểm tra các trường dữ liệu bắt buộc như \texttt{joints3d}, \texttt{pose}, \texttt{trans} và \texttt{shape}, kiểm tra số frame và kích thước của từng trường, rồi chuẩn hóa kiểu dữ liệu, tên trường và đơn vị đo để các mô-đun sau có thể xử lý thống nhất.

\begin{enumerate}[leftmargin=*, label=\arabic*.]
    \item Đọc metadata của hai video, gồm số frame, FPS và kích thước ảnh.
    \item Đọc output monocular tương ứng với từng video \cite{shin2024wham}.
    \item Kiểm tra sự tồn tại của các trường dữ liệu bắt buộc: \texttt{joints3d}, \texttt{pose}, \texttt{shape}, \texttt{trans} \cite{smpl2015}.
    \item Kiểm tra số frame và kích thước của từng trường dữ liệu.
    \item Chuẩn hóa kiểu dữ liệu, tên trường và đơn vị đo.
    \item Trả về hai chuỗi pose thô đã hợp lệ hóa.
\end{enumerate}

Rủi ro chính của mô-đun này là dữ liệu đầu vào không đầy đủ hoặc không đồng nhất giữa hai view. Nếu một số frame thiếu khớp 3D hoặc thiếu tham số SMPL, hệ thống cần đánh dấu để loại khỏi đánh giá hoặc nội suy trong các bước sau. Việc kiểm tra đầu vào giúp tránh lan truyền lỗi định dạng sang các mô-đun alignment và fusion.

\section{Module 2: Đồng bộ thời gian giữa hai chuỗi pose}
\label{sec:temporal_sync}

Mô-đun đồng bộ thời gian nhằm đảm bảo hai chuỗi pose từ hai camera biểu diễn cùng một thời điểm chuyển động. Nếu hai chuỗi chưa được đồng bộ, việc so sánh pose giữa hai góc nhìn có thể tạo ra sai lệch giả, vì hai pose tương ứng thực chất thuộc hai thời điểm khác nhau. Đồng bộ thời gian là điều kiện cần trước khi áp dụng consistency check hoặc fusion đa góc nhìn.

Trong trường hợp hai video có cùng FPS và được ghi đồng thời, đồng bộ có thể thực hiện bằng cách ghép frame cùng chỉ số. Trong trường hợp tồn tại độ lệch thời gian, hệ thống cần ước lượng offset giữa hai chuỗi. Offset có thể được xác định từ timestamp, từ sự kiện đồng bộ ban đầu, hoặc từ tín hiệu chuyển động suy ra từ pose, chẳng hạn quỹ đạo khớp hông, cổ tay, cổ chân hoặc biến thiên góc khớp theo thời gian. Việc sử dụng tín hiệu chuyển động để căn chỉnh chuỗi phù hợp với mục tiêu khai thác temporal information trong 3D human pose estimation \cite{mtftransformer2023,shin2024wham}.

Đầu vào của mô-đun là hai chuỗi pose thô $S^{(1)}$ và $S^{(2)}$. Đầu ra là hai chuỗi pose đã đồng bộ $\bar{S}^{(1)}$ và $\bar{S}^{(2)}$ có cùng độ dài $T$ và có quan hệ frame-to-frame.

\begin{enumerate}[leftmargin=*, label=\arabic*.]
    \item Kiểm tra FPS và số frame của hai video.
    \item Nếu có timestamp, ghép frame theo timestamp gần nhất.
    \item Nếu không có timestamp, ước lượng offset bằng tín hiệu chuyển động từ các khớp ổn định.
    \item Cắt bỏ các frame không có cặp tương ứng.
    \item Trả về hai chuỗi pose có cùng độ dài và tương ứng theo thời gian.
\end{enumerate}

Rủi ro của bước đồng bộ là offset bị ước lượng sai khi chuyển động quá chậm, quá lặp lại hoặc khi một view có nhiều frame lỗi. Vì vậy, kết quả đồng bộ cần được kiểm tra bằng trực quan hóa hoặc bằng sai khác chuyển động giữa hai chuỗi sau khi căn chỉnh.

\section{Module 3: Căn chỉnh pose về không gian so sánh chung}
\label{sec:pose_alignment}

Sau khi đồng bộ thời gian, hai chuỗi pose vẫn có thể nằm trong hai hệ tọa độ khác nhau. Mô-đun căn chỉnh có nhiệm vụ đưa hai chuỗi pose về một không gian so sánh chung trước khi kiểm tra nhất quán đa góc nhìn. Đây là bước cần thiết vì sai khác giữa hai pose chỉ có ý nghĩa khi chúng được biểu diễn trong cùng hệ quy chiếu.

Phương pháp căn chỉnh có thể sử dụng thông số camera nếu có. Trong trường hợp thông số camera không đầy đủ, hệ thống có thể dùng căn chỉnh hình học dựa trên tập khớp tương ứng. Các nghiên cứu multi-view như Gu và cộng sự sử dụng Procrustes alignment trong bài toán ước lượng pose 3D từ nhiều góc nhìn động \cite{gu2024multiview}. Trong đồ án, căn chỉnh được sử dụng ở vai trò hậu xử lý nhằm biến đổi hai chuỗi pose về không gian so sánh, không nhằm giải bài toán camera calibration đầy đủ.

Với mỗi frame $t$, cần tìm phép biến đổi $A_t$ sao cho pose từ view thứ hai gần với pose từ view thứ nhất trên một tập khớp neo $A$. Phép biến đổi có thể gồm quay, tịnh tiến và tỉ lệ:
\[
\tilde{X}^{(2)}_t = s_t R_t X^{(2)}_t + d_t
\]
trong đó $R_t$ là ma trận quay, $s_t$ là hệ số tỉ lệ, và $d_t$ là vector tịnh tiến. Cách mô hình hóa này tương ứng với một phép biến đổi similarity, thường được dùng trong căn chỉnh hình học giữa hai tập điểm 3D \cite{gu2024multiview}. Sau căn chỉnh, hệ thống thu được hai chuỗi $\tilde{X}^{(1)}$ và $\tilde{X}^{(2)}$ trong cùng không gian so sánh.

\begin{enumerate}[leftmargin=*, label=\arabic*.]
    \item Chọn tập khớp neo ổn định, ví dụ pelvis, hip, shoulder và spine.
    \item Với mỗi frame, lấy các khớp tương ứng từ hai view.
    \item Ước lượng phép biến đổi gồm quay, tịnh tiến và tỉ lệ giữa hai tập khớp \cite{gu2024multiview}.
    \item Áp dụng phép biến đổi cho toàn bộ pose của view cần căn chỉnh.
    \item Làm mượt tham số biến đổi theo thời gian nếu cần.
    \item Trả về hai chuỗi pose trong cùng không gian so sánh.
\end{enumerate}

Rủi ro của bước căn chỉnh là phép biến đổi bị ảnh hưởng bởi các khớp lỗi. Để giảm rủi ro này, tập khớp neo nên ưu tiên các khớp trung tâm cơ thể và có thể loại bỏ các khớp có sai khác quá lớn trước khi ước lượng phép biến đổi. Nếu có thông tin camera đáng tin cậy, thông tin này nên được ưu tiên để giảm phụ thuộc vào pose lỗi.

\section{Module 4: Kiểm tra nhất quán đa góc nhìn}
\label{sec:consistency_check}

Mô-đun kiểm tra nhất quán đa góc nhìn có nhiệm vụ phát hiện các khớp hoặc frame có khả năng sai lệch. Nguyên tắc chính là nếu hai camera quan sát cùng một chuyển động, thì hai pose sau căn chỉnh phải tương thích với nhau ở các vùng cơ thể đáng tin cậy. Nguyên tắc này liên quan trực tiếp đến multi-view consistency trong CanonPose, MTF-Transformer và các hướng uncalibrated multi-view pose estimation \cite{canonpose2021,mtftransformer2023,probtriangulation2023}.

Với mỗi frame $t$ và mỗi khớp $j$, sai khác liên góc nhìn được tính bằng:
\[
D_{t,j} = \left\| \tilde{X}^{(1)}_{t,j} - \tilde{X}^{(2)}_{t,j} \right\|_2
\]
Trong đó $\tilde{X}^{(1)}_{t,j}$ và $\tilde{X}^{(2)}_{t,j}$ là tọa độ khớp $j$ tại frame $t$ sau khi hai pose đã được căn chỉnh. Giá trị $D_{t,j}$ càng lớn thì mức độ mâu thuẫn giữa hai góc nhìn tại khớp đó càng cao. Cách đo sai khác theo vị trí khớp có liên hệ với các độ đo tọa độ khớp phổ biến trong 3D pose estimation, nhưng không đủ để đánh giá toàn bộ tính hợp lý của pose \cite{manipose2024}.

Ngoài sai khác vị trí, hệ thống có thể xét thêm sai khác chiều dài xương, sai khác vận tốc và sai khác gia tốc. Cách đánh giá này phù hợp với nhận định của ManiPose rằng các thước đo tọa độ trung bình như MPJPE, P-MPJPE hoặc PCK có thể bỏ sót một số lỗi về tính nhất quán cấu trúc của pose \cite{manipose2024}. Vì vậy, consistency check không nên chỉ dựa vào một độ đo duy nhất.

\begin{table}[H]
\centering
\caption{Các tín hiệu dùng trong kiểm tra nhất quán}
\label{tab:consistency_signals}
\begin{tabular}{|p{3.4cm}|p{4.1cm}|p{3.6cm}|p{3.6cm}|}
\hline
\textbf{Tín hiệu} & \textbf{Cách tính} & \textbf{Dữ liệu cần có} & \textbf{Ứng dụng trong pipeline} \\
\hline
Sai khác vị trí khớp & Khoảng cách Euclidean giữa hai view sau căn chỉnh & Hai chuỗi aligned\_joints3d & Phát hiện mâu thuẫn trực tiếp tại từng khớp \cite{manipose2024} \\
\hline
Sai khác chiều dài xương & Chênh lệch độ dài giữa hai khớp nối với nhau & Skeleton và joints3d & Phát hiện biến dạng skeleton và lỗi cấu trúc \cite{keller2023skin} \\
\hline
Sai khác vận tốc & Chênh lệch chuyển động giữa hai frame liên tiếp & Chuỗi pose theo thời gian & Phát hiện thay đổi bất thường theo thời gian \cite{shin2024wham,physcap2020} \\
\hline
Sai khác gia tốc & Chênh lệch vận tốc theo thời gian & Chuỗi pose theo thời gian & Phát hiện rung giật hoặc nhiễu cục bộ \cite{physcap2020} \\
\hline
Sai khác SMPL/joint & So sánh giữa joint sinh từ SMPL và joint đầu vào & Tham số SMPL và joints & Kiểm tra sự phù hợp giữa biểu diễn khớp và mô hình cơ thể \cite{smpl2015,smplify2016} \\
\hline
Confidence đầu vào & Điểm tin cậy từ mô hình pose nếu có & Output của mô hình monocular & Hỗ trợ chọn view đáng tin cậy hơn khi fusion \cite{mtftransformer2023} \\
\hline
\end{tabular}
\end{table}

Đầu ra của mô-đun là một ma trận nghi ngờ $E \in \{0,1\}^{T \times J}$, trong đó $E_{t,j}=1$ nếu khớp $j$ tại frame $t$ được đánh dấu có khả năng lỗi. Ngoài ra, hệ thống cũng sinh ra ma trận điểm sai khác $D$ và điểm tin cậy ban đầu $C$ để sử dụng trong bước phân xử \cite{vu2026multiperspective}.

\begin{enumerate}[leftmargin=*, label=\arabic*.]
    \item Với mỗi frame $t$ và khớp $j$, tính sai khác vị trí $D_{t,j}$.
    \item Tính thêm sai khác chiều dài xương, vận tốc hoặc gia tốc nếu có dữ liệu \cite{keller2023skin,physcap2020}.
    \item Chuẩn hóa các sai khác về cùng thang đo.
    \item So sánh sai khác với ngưỡng hoặc thống kê theo chuỗi.
    \item Đánh dấu các khớp/frame có sai khác bất thường vào ma trận $E$.
    \item Tạo điểm tin cậy ban đầu $C$ cho từng khớp và từng view.
\end{enumerate}

Rủi ro của bước này là đánh dấu nhầm chuyển động nhanh hợp lệ thành lỗi. Vì vậy, ngưỡng phát hiện không nên chỉ dựa trên sai khác tuyệt đối, mà cần xét thêm ngữ cảnh thời gian và cấu trúc skeleton. Các khớp ở cổ tay hoặc cổ chân có thể chuyển động nhanh hơn khớp trung tâm, nên ngưỡng phát hiện cần được thiết kế theo nhóm khớp.

\section{Module 5: Phân xử và hợp nhất pose}
\label{sec:judgment_fusion}

Sau khi xác định các khớp hoặc frame nghi ngờ, hệ thống cần phân xử giữa hai kết quả monocular. Nếu hai view cho kết quả không thống nhất, việc lấy trung bình trực tiếp có thể làm lan truyền lỗi từ view kém tin cậy sang kết quả cuối cùng. Vì vậy, mô-đun phân xử sử dụng điểm tin cậy cục bộ để quyết định nguồn thông tin nào nên được ưu tiên \cite{vu2026multiperspective}.

Điểm tin cậy của một khớp có thể được tính từ ba nhóm tín hiệu. Nhóm thứ nhất là độ nhất quán liên góc nhìn: khớp có sai khác nhỏ giữa hai view thường đáng tin cậy hơn. Nhóm thứ hai là độ ổn định theo thời gian: khớp có quỹ đạo mượt và không có biến đổi đột ngột thường đáng tin cậy hơn. Nhóm thứ ba là tính hợp lý của skeleton: khớp duy trì quan hệ hợp lý với các khớp lân cận thường được ưu tiên hơn. Các nhóm tín hiệu này tương ứng với multi-view consistency, temporal stability và biomechanical plausibility đã được sử dụng hoặc thảo luận trong nhiều nghiên cứu liên quan \cite{canonpose2021,physcap2020,keller2023skin}.

Giả sử $C^{(1)}_{t,j}$ và $C^{(2)}_{t,j}$ là điểm tin cậy của khớp $j$ tại frame $t$ từ hai view. Pose hợp nhất có thể được tính theo trung bình có trọng số:
\[
X^{F}_{t,j}
=
\frac{
C^{(1)}_{t,j}\tilde{X}^{(1)}_{t,j}
+
C^{(2)}_{t,j}\tilde{X}^{(2)}_{t,j}
}{
C^{(1)}_{t,j}+C^{(2)}_{t,j}+\epsilon
}
\]
trong đó $\epsilon$ là hằng số nhỏ để tránh chia cho 0. Cách hợp nhất theo trọng số tin cậy có liên hệ với các cơ chế fusion sử dụng confidence trong các hệ thống multi-view pose estimation \cite{mtftransformer2023}. Nếu một trong hai view có độ tin cậy thấp rõ rệt, hệ thống có thể chọn trực tiếp view còn lại thay vì lấy trung bình.

\begin{enumerate}[leftmargin=*, label=\arabic*.]
    \item Với mỗi frame $t$ và khớp $j$, lấy điểm tin cậy từ hai view.
    \item Nếu cả hai view có độ tin cậy cao, hợp nhất bằng trung bình có trọng số \cite{mtftransformer2023}.
    \item Nếu một view có độ tin cậy thấp, ưu tiên view còn lại.
    \item Nếu cả hai view đều không tin cậy, nội suy từ các frame lân cận hoặc dùng ràng buộc skeleton \cite{keller2023skin}.
    \item Kiểm tra lại chiều dài xương và tính liên tục sau hợp nhất.
    \item Trả về chuỗi pose hợp nhất $X^F$.
\end{enumerate}

Rủi ro của mô-đun phân xử là điểm tin cậy được tính sai, đặc biệt khi cả hai view cùng sai tại một vùng cơ thể. Trong trường hợp này, hệ thống cần dựa thêm vào ràng buộc thời gian và ràng buộc động học thay vì chỉ dựa vào so sánh hai view. Đây là lý do bước refinement cuối cùng vẫn cần thiết sau khi đã fusion.

\section{Module 6: Hiệu chỉnh pose bằng ràng buộc thời gian và động học}
\label{sec:pose_refinement}

Mô-đun hiệu chỉnh pose nhận chuỗi pose hợp nhất $X^F$ và tạo ra chuỗi pose cuối cùng $\hat{X}$. Mục tiêu của bước này là giảm rung giật theo thời gian, duy trì chiều dài xương hợp lý và tăng sự phù hợp với mô hình cơ thể. Các nghiên cứu như PhysCap, SKEL và OpenSim cho thấy chuyển động người cần được đánh giá không chỉ bằng tọa độ khớp, mà còn bằng tính hợp lý vật lý, cấu trúc xương và ràng buộc cơ sinh học \cite{physcap2020,keller2023skin,delp2007opensim}.

Ràng buộc thời gian có thể được mô tả bằng sai khác vận tốc hoặc gia tốc giữa các frame liên tiếp:
\[
L_{\text{temp}} =
\sum_{t=2}^{T-1}
\left\|
X_{t+1} - 2X_t + X_{t-1}
\right\|_2^2
\]
Thành phần này hạn chế các biến đổi đột ngột và giảm jitter, một vấn đề thường được thảo luận trong các hướng monocular motion capture và physics-based refinement \cite{shin2024wham,physcap2020}. Tuy nhiên, hệ số của ràng buộc thời gian cần được lựa chọn thận trọng để không làm mất các chuyển động nhanh có thật.

Ràng buộc chiều dài xương được dùng để duy trì cấu trúc skeleton:
\[
L_{\text{bone}} =
\sum_{t=1}^{T}
\sum_{(i,j)\in B}
\left(
\|X_{t,i}-X_{t,j}\|_2 - l_{i,j}
\right)^2
\]
trong đó $B$ là tập các cặp khớp tạo thành xương và $l_{i,j}$ là chiều dài tham chiếu của xương $(i,j)$. Chiều dài tham chiếu có thể được lấy từ frame tin cậy, từ trung vị theo chuỗi hoặc từ mô hình cơ thể. Ràng buộc này phù hợp với các nghiên cứu nhấn mạnh vai trò của skeleton và biomechanical plausibility trong biểu diễn cơ thể người \cite{keller2023skin,delp2007opensim}.

Nếu hệ thống sử dụng biểu diễn SMPL, bước refinement có thể cập nhật tham số pose, shape hoặc translation rồi sinh lại joints và vertices. SMPLify sử dụng tối ưu hóa để khớp mô hình SMPL với quan sát 2D \cite{smplify2016}. Learnable SMPLify là một hướng tham khảo vì sử dụng mô hình học sâu để giải bài toán inverse kinematics với chi phí thấp hơn tối ưu lặp truyền thống \cite{yang2025learnable}. Trong đồ án, thành phần này được xem như bước hiệu chỉnh cuối, không phải mô hình ước lượng pose ban đầu.

\begin{enumerate}[leftmargin=*, label=\arabic*.]
    \item Nhận pose hợp nhất từ mô-đun phân xử.
    \item Tính ràng buộc thời gian dựa trên vận tốc hoặc gia tốc \cite{physcap2020}.
    \item Tính ràng buộc chiều dài xương dựa trên skeleton tham chiếu \cite{keller2023skin}.
    \item Nếu dùng SMPL, cập nhật tham số pose, shape hoặc trans tương ứng \cite{smpl2015,smplify2016}.
    \item Sinh lại joints3D và vertices từ tham số đã hiệu chỉnh.
    \item Trả về chuỗi pose cuối cùng để đánh giá và trực quan hóa.
\end{enumerate}

Rủi ro của bước refinement là làm quá mượt chuyển động hoặc ép pose về một cấu trúc không phản ánh đúng động tác thực. Vì vậy, trong thực nghiệm cần đánh giá đồng thời sai số tọa độ, độ ổn định thời gian và các chỉ số cấu trúc như biến thiên chiều dài xương. Cách đánh giá đa tiêu chí này phù hợp với nhận định của ManiPose rằng các chỉ số chuẩn như MPJPE, P-MPJPE và PCK chưa đủ để phản ánh toàn bộ tính nhất quán của pose \cite{manipose2024}.

\section{Đánh giá phương pháp đề xuất}
\label{sec:evaluation_design}

Việc đánh giá phương pháp cần được thực hiện trên cả kết quả cuối cùng và các kết quả trung gian sau từng mô-đun. Cách đánh giá theo từng bước giúp xác định thành phần nào đóng góp vào thay đổi của chất lượng pose, đồng thời tránh kết luận chung chung về toàn bộ pipeline. Các cấu hình ablation và bộ tiêu chí đánh giá được chuyển sang Chương 5 để gắn trực tiếp với phần thực nghiệm.

Ngoài đánh giá định lượng, trực quan hóa vẫn cần thiết để kiểm tra các lỗi khó phản ánh bằng chỉ số trung bình. Hệ thống có thể trực quan hóa skeleton 3D, mesh SMPL và keypoints chiếu ngược lên video 2D. Trực quan hóa skeleton và mesh phù hợp với các biểu diễn được sử dụng trong SMPL, SMPLify và các hệ thống motion capture dựa trên mô hình cơ thể \cite{smpl2015,smplify2016}. Các kết quả trực quan này hỗ trợ phát hiện lỗi cục bộ như đảo hướng chi, rung giật ở cổ tay hoặc cổ chân, sai lệch chiều sâu và pose không tự nhiên.

\section{Kết chương}
\label{sec:method_conclusion}

Chương này đã trình bày phương pháp đề xuất cho bài toán phát hiện và hiệu chỉnh lỗi ước lượng tư thế người 3D. Pipeline được thiết kế theo hướng hậu xử lý hai chuỗi pose monocular, gồm các mô-đun nạp dữ liệu, đồng bộ thời gian, căn chỉnh pose, kiểm tra nhất quán đa góc nhìn, phân xử hợp nhất và hiệu chỉnh bằng ràng buộc thời gian, ràng buộc động học và biểu diễn SMPL.

Các mô-đun trong pipeline được mô tả theo định dạng input, output và quy trình xử lý cụ thể. Chương cũng đã trình bày các rủi ro chính của từng bước và thiết kế đánh giá dự kiến để kiểm chứng vai trò của từng thành phần. Chương tiếp theo sẽ trình bày chi tiết cơ sở triển khai, hàm mục tiêu hoặc các thành phần kỹ thuật được sử dụng trong quá trình cài đặt phương pháp đề xuất.
\end{document}
